\uuid{Xcce}
\titre{Coordonnées}
\auteur{bodin ; hanani ; mzari}
\organisation{exo7}
\licence{CC BY-NC-SA 4.0}
\datecreate{2020-01-20}
\niveau{L1}
\module{Géométrie}
\chapitre{Géométrie du plan}
\sousChapitre{Distance, norme, produit scalaire}
\difficulte{1}
\theme{coordonnées}
\exo7qcmid{lille-1.geometrie-plane.004}

\contenu{
\begin{question}
\qclassification{140.01,140.02}
\qtags{motcle=coordonnées}
Dans un repère orthonormé direct, on considère le point $A$ de coordonnées polaires $r=2$ et $\displaystyle \theta =\frac{\pi}{6}$. Quelles sont les coordonnées cartésiennes $(x,y)$ de $A$ ?
\begin{answers}
    \bad{$x=2$ et $y=2$}
    \good{$x=\sqrt{3}$ et $y=1$}
    \bad{$x=1$ et $y=\sqrt{3}$}
    \bad{$x=1$ et $y=1$}
\end{answers}
\begin{explanations}
Les deux systèmes de coordonnées sont reliés par les relations $x=r\cos \theta$ et $y=r\sin \theta $.
\end{explanations}
\end{question}
}
